\chapter{Alkalmazott technológiák}
\label{ch:technologiak}
% Terjedelem: kb. 8–10 oldal (~14 000 karakter)
% Nem csak felsorolás: mindegyiknél MIÉRT ezt választottad, mi az alternatíva.

\section{Webes architektúra és a REST elv}
% Kliens–szerver, REST megszorításai \cite{fielding2000architectural}, HTTP metódusok.
\todo{Szöveg}

\section{A backend technológiái}
\subsection{FastAPI és az ASGI}
% Típusannotációk, Pydantic validáció, automatikus OpenAPI/Swagger \cite{fastapi_docs}.
\todo{Szöveg}

\subsection{Adatbázis: PostgreSQL és SQLAlchemy ORM}
% Relációs modell, ORM előnyei \cite{sqlalchemy_docs}.
\todo{Szöveg}

\subsection{Hitelesítés: JWT és bcrypt}
% Állapotmentes token \cite{rfc7519}, jelszó-hash sózással \cite{provos1999bcrypt}.
\todo{Szöveg}

\section{Valós idejű kommunikáció: WebSocket és Socket.IO}
% \cite{rfc6455,socketio_docs}
\todo{Szöveg}

\section{A frontend technológiái: React, Next.js, Tailwind CSS}
% Komponensalapú UI, App Router, kliensoldali állapotkezelés \cite{nextjs_docs}.
\todo{Szöveg}

\section{Konténerizáció és sandboxing: Docker}
% Konténer vs. virtuális gép, névterek (namespaces), cgroups, képességek (capabilities)
% \cite{merkel2014docker,docker_docs}. Nem megbízható kód futtatásának kockázatai \cite{owasp_top10}.
\todo{Szöveg}
